\documentclass[11pt]{article}

\usepackage{amsmath, amssymb}
\usepackage{geometry}
\usepackage{booktabs}
\usepackage{hyperref}
\usepackage{parskip}
\usepackage{titlesec}
\usepackage{array}
\usepackage{xcolor}

\geometry{margin=1.25in}

\title{Adding a Slow Recovery Variable to the Yuval Neuron Model\\
\large C.\ elegans Locomotion Circuit}
\author{Distribution of Oscillators Project}
\date{\today}

\begin{document}

\maketitle

\begin{abstract}
The piecewise-linear Yuval neuron model is intrinsically bistable: with zero applied current,
the plateau $V = H$ is a stable fixed point and the neuron cannot return to rest without
an external input. This document motivates and derives a minimal extension --- a slow
intrinsic recovery variable $w$ --- that converts the bistable system into a limit-cycle
oscillator, eliminating plateau locking while preserving the qualitative dynamics required
for the Distribution of Oscillators framework.
\end{abstract}

% =============================================================================
\section{The Problem: Structural Bistability}

Each neuron in the current model satisfies

\begin{equation}
    C\,\dot{V} = g\,f(V) - I_{\mathrm{syn}} + I_{\mathrm{app}},
    \label{eq:base}
\end{equation}

where $f(V)$ is the piecewise-linear membrane nonlinearity:

\begin{equation}
f(V) = \begin{cases}
    m_1(-V + L)         & V < L \\
    m_2(V - L)(V - T)   & L \leq V \leq T \\
    m_3(V - H)(V - T)   & T < V \leq H \\
    m_4(-V + H)         & V > H.
\end{cases}
\label{eq:fv}
\end{equation}

By construction, $f(L) = 0$ and $f(H) = 0$ for any choice of parameters. When
$I_{\mathrm{app}} = 0$ and $I_{\mathrm{syn}} = 0$, the fixed points of
\eqref{eq:base} are exactly the zeros of $f(V)$: rest ($V = L$), threshold ($V = T$,
unstable), and plateau ($V = H$). Both $L$ and $H$ are stable, so the neuron is
\emph{bistable by construction}.

Synaptic fatigue (the $\mathrm{SF}$ variable) addresses this in a \emph{network} context:
a fatiguing inhibitory partner can pull the active neuron off the plateau. However, an
isolated neuron receiving $I_{\mathrm{app}} = 0$ has no such partner and will remain
locked at the plateau indefinitely.

Attempts to resolve this by pushing $T$ close to $H$ shrink the basin of attraction of the
plateau but do not eliminate it, because $f(H) = 0$ is a structural property of the
piecewise-linear form, not a consequence of any particular parameter choice.

% =============================================================================
\section{The Fix: A Slow Recovery Variable}

The standard approach in computational neuroscience is to augment the voltage equation with
a slow variable $w$ that plays the role of an intrinsic recovery current. The extended system is

\begin{align}
    C\,\dot{V} &= g\,f(V) - w - I_{\mathrm{syn}} + I_{\mathrm{app}},
    \label{eq:V}\\
    \tau_w\,\dot{w} &= w_\infty(V) - w,
    \label{eq:w}
\end{align}

where $\tau_w \gg \tau_V$ (slow timescale) and $w_\infty(V)$ is a monotonically increasing
target function. Three natural choices for $w_\infty(V)$ are discussed in
Section~\ref{sec:choices}.

\subsection{Geometric Interpretation}

The $V$-nullcline ($\dot{V} = 0$) is now

\begin{equation}
    w = g\,f(V) + I_{\mathrm{app}},
    \label{eq:Vnull}
\end{equation}

which is an $N$-shaped curve inherited from $f(V)$. The $w$-nullcline ($\dot{w} = 0$) is

\begin{equation}
    w = w_\infty(V),
    \label{eq:wnull}
\end{equation}

a monotonically increasing curve. The system has a fixed point wherever the two nullclines
intersect.

When $I_{\mathrm{app}} = 0$, the $V$-nullcline passes through $(L,\,0)$ and $(H,\,0)$.
If $w_\infty$ is chosen so that its curve intersects the $V$-nullcline \emph{only on the
left (rest) branch}, then the plateau fixed point is destroyed and the only attractor is
rest. For larger $I_{\mathrm{app}}$, the $V$-nullcline shifts upward, the intersection
can move to the right branch, and the system either settles at a depolarized fixed point
or --- if the intersection is in the negative-slope region of $f(V)$ --- executes a stable
limit cycle (oscillation).

This is exactly the FitzHugh--Nagumo geometry applied to the piecewise-linear model.

% =============================================================================
\section{Choices for \texorpdfstring{$w_\infty(V)$}{w\_inf(V)}}
\label{sec:choices}

\subsection{Linear (simplest)}

\begin{equation}
    w_\infty(V) = \alpha\,(V - L),
\end{equation}

with $\alpha > 0$. At rest $V = L$, so $w_\infty = 0$ and $w \to 0$: the recovery
variable vanishes at rest. At the plateau $V = H$,
$w_\infty = \alpha(H - L) > 0$, providing a net hyperpolarizing current that
pulls the neuron off the plateau. This is the most transparent choice and
introduces a single free parameter $\alpha$.

The condition for the plateau fixed point to be eliminated is that the $w$-nullcline
evaluated at $H$ lies \emph{above} the $V$-nullcline at $H$ (which equals zero):

\begin{equation}
    \alpha(H - L) > 0,
\end{equation}

which is always satisfied for $\alpha > 0$. Any positive $\alpha$ eliminates the
plateau fixed point when $I_{\mathrm{app}} = 0$.

\subsection{Sigmoid (biophysically motivated)}

\begin{equation}
    w_\infty(V) = \frac{w_{\max}}{1 + e^{-k_w(V - V_{w})}}
\end{equation}

A sigmoid centered near $T$ or above mimics slow K$^+$ channel activation. This is
the biophysically natural form but introduces three parameters ($w_{\max}$, $k_w$,
$V_w$), most of which are not constrained by existing data for this model.

\subsection{Threshold-gated (piecewise, consistent with existing style)}

\begin{equation}
    w_\infty(V) = \begin{cases} 0 & V \leq T \\ \beta(V - T) & V > T, \end{cases}
\end{equation}

Recovery is zero below threshold and grows linearly above it. This is consistent
with the piecewise-linear style of the rest of the model and keeps the number of
free parameters minimal (one: $\beta$). At rest $w_\infty(L) = 0$; at the plateau
$w_\infty(H) = \beta(H - T) > 0$.

% =============================================================================
\section{Timescale Separation and Oscillation}

For the extended system to produce oscillations, the recovery variable must be slow
enough that the voltage can traverse the active phase before $w$ equilibrates, but
fast enough that $w$ accumulates meaningfully during a plateau. Concretely:

\begin{itemize}
    \item $\tau_w \ll$ plateau duration: $w$ builds up during the plateau and
          terminates it.
    \item $\tau_w \gg$ spike upswing: the upswing is fast enough that $w$ does not
          prevent activation.
\end{itemize}

A rough estimate: plateau duration is set by the synaptic fatigue timescale
($1/a \approx 1/0.000035 \approx 28{,}000$~ms per unit drive), while the
membrane time constant is $C/g = 7$~ms. A reasonable starting range for
$\tau_w$ is $100$--$1000$~ms, which is slow relative to the membrane but fast
relative to the network rhythm.

% =============================================================================
\section{Interaction with Synaptic Fatigue}

The existing synaptic fatigue variable $\mathrm{SF}$ acts as a slow recovery
variable at the \emph{network} level: it terminates inhibition from the active
neuron, allowing its partner to escape. The new $w$ variable acts at the
\emph{cell} level: it terminates the plateau of the active neuron directly.

These two mechanisms are complementary:

\begin{center}
\begin{tabular}{lll}
\toprule
Mechanism       & Acts on          & Timescale \\
\midrule
Synaptic fatigue $\mathrm{SF}$ & Inhibitory current onto partner & $1/a$, $1/b$ (very slow) \\
Recovery variable $w$          & Active neuron's own plateau     & $\tau_w$ (tunable) \\
\bottomrule
\end{tabular}
\end{center}

With both mechanisms present, the half-center oscillator has two routes to burst
termination. This makes the oscillation more robust and reduces dependence on the
partner neuron's fatigue — a single isolated neuron can now oscillate intrinsically
when $I_{\mathrm{app}} > 0$, and return to rest when $I_{\mathrm{app}} = 0$.

% =============================================================================
\section{Implementation in the Existing Codebase}

The minimal change to \texttt{src/functions.py} is to add the $w$-dynamics:

\begin{equation}
    \dot{w} = \frac{w_\infty(V) - w}{\tau_w},
\end{equation}

and to modify \texttt{src/network.py} to expand the state vector from
$[V_1, \ldots, V_N, \mathrm{SF}_1, \ldots, \mathrm{SF}_N]$ to
$[V_1, \ldots, V_N, \mathrm{SF}_1, \ldots, \mathrm{SF}_N, w_1, \ldots, w_N]$,
subtracting $w_i$ from the $\dot{V}_i$ equation.

New global parameters to add to \texttt{src/settings.py}:
\begin{itemize}
    \item \texttt{tau\_w}: recovery timescale (ms), suggested initial value $200$~ms.
    \item \texttt{w\_scale}: amplitude scaling ($\alpha$ or $\beta$ above), suggested
          initial value $0.01$~nS/mV.
    \item \texttt{V\_w}: threshold for the piecewise variant (= $T$ in the simplest case).
\end{itemize}

% =============================================================================
\section{Summary}

The bistability of the Yuval model is structural: $f(H) = 0$ by construction, making
the plateau a permanent fixed point when $I_{\mathrm{app}} = 0$. A slow recovery
variable $w$ resolves this by adding a nullcline that intersects $f(V)$ only on the
rest branch at zero drive. The piecewise-linear form $w_\infty(V) = \beta\max(V-T,\,0)$
is the most natural extension consistent with the existing model style, introduces one
free parameter, and ensures the neuron returns to rest under zero drive while remaining
capable of oscillation under sufficient excitation.

\end{document}
